% Lösungen Ableitung und Änderungsrate · Die Rate als Funktion · Prüfungs-Fokus Abitur GK (zusammenbau v0.6)
\einheitenkopf{Einheit 4 · Die Rate als Funktion}

\erg{1}{a) Bestand (Höhe, Anzahl, Menge) – größte Rate bei den Nullstellen der zweiten Ableitung: die Zuflussrate ist $f'$, ihr Maximum liegt bei $f'' = 0$ \quad b) $r'(t) = 0$ ⇔ $t = 6$ (der e-Faktor hat keine Nullstelle); nach 6 Stunden ist die Zuflussrate am größten \quad c) $T'(t) = -0{,}3t^2 + 3{,}6t$, $T''(t) = -0{,}6t + 3{,}6 = 0$ ⇔ $t = 6$; $T'(6) = 10{,}8$ °C je h}
\erg{2}{a) $f'(x) = 1{,}5x^2 - 3$, $f''(x) = 3x = 0$ ⇔ $x = 0$; kleinste Steigung $f'(0) = -3$ (nicht das Minimum von $f$ suchen) \quad b) $p''(x) = (0{,}32x - 1{,}44) \cdot \mathrm{e}^{-0{,}4x} = 0$ ⇔ $x = 4{,}5$; dort ist $p'(4{,}5) \approx -0{,}331 < -0{,}33$ (die mittlere Steigung im Intervall, etwa $-0{,}26$, genügt nicht als Beleg) \quad c) $w''(t) = -6t + 24 = 0$ ⇔ $t = 4$; $w'(4) = 48$ cm je Woche (nicht $t = 4$ als Ergebnis angeben) \quad d) $k''(x) = 0{,}3 \cdot (1 - 0{,}1x) \cdot \mathrm{e}^{-0{,}1x} = 0$ ⇔ $x = 10$; $k'(10) = 3 \cdot \mathrm{e}^{-1} \approx 1{,}10$ Mio. t je Jahr (nicht $k'(x) = 0$ lösen) \quad e) $h'(t) = 23{,}4$ ⇔ $t^2 - 20t + 36 = 0$ ⇔ $t = 2$ oder $t = 18$; $t = 18$ liegt außerhalb der Phase; $h'(0) = 45$: die Geschwindigkeit ist von $t = 0$ bis $t = 2$ mindestens $23{,}4$ m je min, der Zeitraum ist $2$ Minuten lang \quad f) $f'(t) = 9$ ⇔ $t^2 - 8t + 12 = 0$ ⇔ $t = 2$ oder $t = 6$; beide in der Phase: von $t = 2$ bis $t = 6$, die Länge des Zeitraums ist $= 4$ Minuten}
\erg{3}{a) Differenz der Raten $d(t) = a'(t) - b'(t) = 36t - 3t^2 - 6t = 30t - 3t^2$; $d'(t) = 30 - 6t = 0$ ⇔ $t = 5$; $d(5) = 75$ cm je Jahr; Ränder: $d(0) = 0$, $d(12) = -72$ – der Unterschied ist nach 5 Jahren mit $75$ cm je Jahr am größten (nicht die Differenz der Höhen untersuchen) \quad b) $d(t) = s'(t) - u'(t) = 0{,}15t^2 + 2 - 3t$; $d'(t) = 0{,}3t - 3 = 0$ ⇔ $t = 10$; $d(10) = -13$, Ränder $d(0) = 2$, $d(15) = -9{,}25$ – nach 10 Sekunden ist der Unterschied mit $13$ m je s am größten (das zweite Fahrzeug ist schneller) \quad c) $d(t) = a'(t) - b'(t) = 3t^2 - 18t + 15$; $d'(t) = 6t - 18 = 0$ ⇔ $t = 3$ mit $d(3) = -12$; Ränder: $d(0) = 15$, $d(6) = 15$ – der Unterschied ist zu Beginn und nach 6 Monaten mit $15$ Tausend je Monat am größten, die innere Extremstelle liefert nur $12$ \quad d) $T(x) - 6 = 50 \cdot \mathrm{e}^{-0{,}02x} = -50 \cdot (-\mathrm{e}^{-0{,}02x}) = -50 \cdot T'(x)$, also $c = -50$ – die Aussage ist wahr; $c$ ist negativ, weil die Temperatur sinkt \quad e) $m(t) = 80 \cdot \mathrm{e}^{-0{,}25t} = -4 \cdot (-20 \cdot \mathrm{e}^{-0{,}25t}) = -4 \cdot m'(t)$, also $c = -4$ – wahr; der Faktor ist negativ, weil die Menge abnimmt \quad f) $L(t) - 100 = -60 \cdot \mathrm{e}^{-0{,}05t} = -20 \cdot (3 \cdot \mathrm{e}^{-0{,}05t}) = -20 \cdot L'(t)$, also $c = -20$ – wahr; der Abstand zur Vollladung ist proportional zur Laderate \quad g) $B(t) - 200 = 500 \cdot \mathrm{e}^{0{,}04t} = 25 \cdot (20 \cdot \mathrm{e}^{0{,}04t}) = 25 \cdot B'(t)$, also $c = 25$ – wahr; hier ist $c$ positiv, weil die Zahl wächst}
